\subsection{Information disclosure and its relationship to masked kernel coordinates}\label{sec:dpr:privacy}

The tensor compiler above is \emph{not}
zero knowledge. Its first prover
message is the collection of all
partial evaluations $s_v^*$,
represented as an element of $A$.
This is precisely the leakage
surface addressed by later work
\cite{Kochergin2026,FrigoShelat2026}.

\begin{lemma}[The first message reveals every partial evaluation]\label{lem:dpr:leaks-partials}
In an honest execution,
the verifier can calculate
\[
 \boxed{s_v^*=\operatorname{Col}(S_0)_v}
\]
for every $v\in I_\ell$
using only its first received
ring-switching message.
\end{lemma}
\begin{proof}
For an honest prover $S_0=S$.
Theorem~\ref{thm:dpr:column-evaluation}
identifies the $v$-th
tensor column of $S$
with the desired partial
evaluation $s_v^*$.
The column operation is
explicitly available from
the public tensor-field basis,
so the verifier can evaluate it.
\end{proof}

\begin{theorem}[Perfect two-witness distinction from the first message]\label{thm:dpr:not-zk}
For every $\ell\geq1$ and $h\geq1$,
there exist two distinct Boolean
witness tables $w^{(0)},w^{(1)}$
and a public evaluation point
$(r_{\rm lo},r_{\rm hi})$
such that both tables have the same
claimed evaluation $y_{\rm pub}=0$,
but the distributions of the honest
first tensor messages $S_0$
have total variation distance $1$.
Thus the unmasked tensor protocol
does not satisfy statistical
zero knowledge in the
evaluation-only public-input game
of Chapter~3.
\end{theorem}
\begin{proof}
Choose $r_{\rm lo}=0_\ell$
and $r_{\rm hi}=0_h$, both
Boolean evaluation points.
Let $w^{(0)}$ be identically zero.
Choose a nonzero Boolean index
$v_*\in I_\ell$ and let
$w^{(1)}$ be zero everywhere
except at $(v_*,0_h)$,
where it equals $1$.
At the Boolean point
$(0_\ell,0_h)$,
both multilinear extensions
evaluate to the witness entry
$w^{(i)}_{0_\ell,0_h}=0$,
so both public claims agree.

For $w^{(0)}$, every partial
evaluation $s_v^*$ is zero.
For $w^{(1)}$,
the high-axis weights at
$r_{\rm hi}=0_h$ equal
$\operatorname{eq}_h(j,0_h)
=\delta_{j,0_h}$.
Therefore the partial vector
has entry $s_{v_*}^*=1$
and every other entry zero.
By Theorem~\ref{thm:dpr:column-evaluation},
the honest first messages have
different $v_*$-th columns:
$0$ and $1$, respectively.
The deterministic test asking
whether the $v_*$-th column
equals zero outputs $1$
with probability $1$ for the
first witness and $0$ for
the second.
Lemma~\ref{lem:security:events}
therefore makes their total
variation distance equal to $1$.
\end{proof}

\begin{definition}[Masked field-column interface]\label{def:dpr:masked-column-interface}
To avoid exposing $s_v^*$ directly,
the field-affine interface of
Chapter~7 uses the public vector
\[
 m_v:=s_v^*+p_v
\]
for hidden independent uniform
pads $p_v\in L$. It also introduces
the hidden transposed correction
$c^p=\mathcal T(p)$
and checks the contraction
$D_c(\theta)=\Delta_p(\theta)$
using the trace-coordinate
checker of Chapter~6.
\end{definition}

\begin{lemma}[The basic masked column view does not distinguish]\label{lem:dpr:masked-column-uniform}
For any two fixed Boolean witness
tables, the distributions of their
basic public masked vectors $m$
from Definition~\ref{def:dpr:masked-column-interface}
are identical and uniform on $L^I$.
\end{lemma}
\begin{proof}
For either fixed witness,
the partial vector $s^*$
is fixed. Lemma~\ref{lem:pack:masked-uniform}
proves that independent uniform
$p\in L^I$ makes $s^*+p$
uniform on $L^I$.
This distribution is the same
for both witnesses.
\end{proof}

\begin{lemma}[Revealing the full hidden correction reverses the mask]\label{lem:dpr:masked-correction-reveal}
If the verifier obtains both
$m=s^*+p$ and the full
transposed pad vector
$c^p=\mathcal T(p)$,
then it can reconstruct
all of $s^*$ exactly.
\end{lemma}
\begin{proof}
Theorem~\ref{thm:pack:transpose-involutive}
gives $p=\mathcal T(c^p)$.
Consequently
$s^*=m+p=m+\mathcal T(c^p)$,
as already proved in
Theorem~\ref{thm:pack:reveal-leak}.
\end{proof}

\begin{proposition}[The two protocol routes are distinct]\label{prop:dpr:routes-distinct}
The original unmasked compiler of
Definition~\ref{def:dpr:full-protocol}
and the masked algebraic checker
of Chapters~6--7 have different
verifier interfaces:
\begin{enumerate}[label=\textup{(\roman*)}]
\item The former sends an
$A$-valued partial evaluation
$S_0$ and executes an
$h$-round degree-two
sumcheck over $A$ with
challenges in $\varphi_1(L)$.
\item The latter exposes
uniform $m\in L^I$,
keeps $p,c$ hidden,
and uses a random kernel
contraction and an anchored
trace-zero consistency checker
over $L$ with challenges
$(\theta,z)$.
\end{enumerate}
Theorem~\ref{thm:dpr:compose-sound}
does not automatically imply
soundness of the latter;
Theorem~\ref{thm:pack:algebraic-sound}
does not automatically imply
zero knowledge or proximity
soundness of the former.
\end{proposition}
\begin{proof}
The first distinction follows
directly by comparing the
typed messages and checks in
Definitions~\ref{def:dpr:full-protocol}
and~\ref{def:dpr:masked-column-interface}.
The two soundness theorems
quantify different acceptance
events and impose different
timing and underlying-oracle
hypotheses. Neither theorem's
hypotheses entail the other's,
so no implication is asserted
without a new reduction proof.
\end{proof}

\begin{remark}[Prior-work attribution]
Diamond--Posen introduce the
tensor-algebra ring-switching
compiler~\cite{DiamondPosen2024}.
Diamond's zero-knowledge
large-field PCS note explicitly
leaves ring-switching outside
its scope~\cite{Diamond2025ZK}.
Kochergin and Frigo--shelat
subsequently give masking
approaches for binary-field
ring switching
\cite{Kochergin2026,FrigoShelat2026}.
Our use of the Boolean kernel
basis as a representation of the
correction has been established
algebraically; a claim that the
complete cryptographic construction
is new or secure requires a
separate comparison and a full
zero-knowledge simulator.
\end{remark}
